\section{\methodname}
\label{sec:method}

Figure~\ref{fig:task-method-overview} summarizes \methodshort{}: construct
dynamic source-free supervision (\S\ref{subsec:counterfactual-pairs}),
replace only the factual source-slot identity relative to Matched
(\S\ref{subsec:source-slot-randomization}), and alternate counterfactual
distillation with preservation updates (\S\ref{subsec:training-objectives}).

\subsection{Structured Dynamic Counterfactual Pairs}
\label{subsec:counterfactual-pairs}

\paragraph{Structured factual bindings.}
A target noun or spatial mask does not specify the event whose consequence
must also be removed. For each mechanism \(m\), we therefore build factual
prompts from registered event fields rather than editing free-form strings.
For training binding \(i\), let
\begin{equation}
  p_i^{\mathrm{f}}
  =
  T_{m_i}^{\mathrm f}(s_i,a_i,r_i,f_i,c_i;w_i),
  \label{eq:factual-prompt}
\end{equation}
where \(s_i\) is the source, \(a_i\) the trigger or interaction, \(r_i\)
the receiver, \(f_i\) the downstream footprint, \(c_i\) the remaining scene
context, and \(w_i\) the registered direct or natural surface realization.
The factual template \(T_{m_i}^{\mathrm f}\) realizes these fields while
preserving their roles.

\paragraph{Dynamic source-free targets.}
Removing the event fields from text alone still does not provide a plausible
dynamic training target. We construct a source-free prompt
\begin{equation}
  p_i^{\mathrm{cf}}
  =
  T_{m_i}^{\mathrm{cf}}(r_i,h_i,c_i;w_i)
  \quad\text{and generate a proxy}\quad
  x_i^{\mathrm{cf}}
  \sim G_{\theta}(\,\cdot\mid p_i^{\mathrm{cf}},n_i),
  \label{eq:distributional-counterfactual}
\end{equation}
where \(h_i\) is the registered no-source receiver state and \(n_i\) is a
generation-only exclusion condition naming forbidden source, trigger, and
footprint content. Using the frozen base generator \(G_{\theta}\), the target
prompt requests that the receiver follow \(h_i\) with natural temporal
variation, while the exclusion condition discourages event leakage. The
generated video is an independently sampled, screened proxy for the ideal
source-free distribution, not the frame-aligned counterfactual of one
realized video. The embedding of \(p_i^{\mathrm{cf}}\) later conditions the
online teacher; \(n_i\) is used for target-media generation only and is not
part of that teacher condition.

\paragraph{Screening and preservation.}
Treatment-blind screening freezes exactly 178 of 192 candidate erasure
bindings per mechanism before either arm trains. Deterministic rejection,
ranking, and tie-breaking govern selection; a version with fewer eligible
targets is invalid. These screened proxies may still contain partial or
uncertain cues. All 18 Wan method/control runs share 36 generic preservation
bindings. Target latents and prompt embeddings are cached and hash-bound
before optimization; the frozen teacher prediction is computed online on the
student's noised target latent. Templates, screening rules, preservation
prompts, and cache contracts appear in \appref{app:counterfactual-data}.

\subsection{Randomizing the Causal-Source Slot}
\label{subsec:source-slot-randomization}

\paragraph{Identity is a shortcut.}
With only a small set of repeated training sources, an adapter can associate
the objective with particular nouns or recurring surface forms. Such a
shortcut can suppress familiar sources without learning a rule that transfers
to a new entity occupying the same event role. We target this ambiguity
directly rather than changing the counterfactual videos, loss weights, or
optimization schedule.

\paragraph{A single source-slot intervention.}
For mechanism \(m\), let \(\mathcal{Q}_m\) be its frozen 64-item source bank,
and let \(q_i\in\mathcal{Q}_{m_i}\) be the deterministically assigned source
for erasure binding \(i\), with \(q_i\neq s_i\). The augmented factual prompt
is rebuilt from the same structured fields:
\begin{equation}
  p_i^{\mathrm{aug}}
  =
  T_{m_i}^{\mathrm f}(q_i,a_i,r_i,f_i,c_i;w_i).
  \label{eq:source-slot-prompt}
\end{equation}
This reconstruction changes the occupant of the source slot; it is not
substring substitution. The committed assignment balances the bank over the
active erasure schedule and is fixed before training rather than sampled
online.

The complete frozen pair banks can be written as
\begin{align}
  \mathcal{B}_{\mathrm{Matched}}
    &= \{(p_i^{\mathrm{f}},x_i^{\mathrm{cf}})\}_{i=1}^{178},
      \label{eq:matched-pairs}\\
  \mathcal{B}_{\mathrm{SRCD}}
    &= \{(p_i^{\mathrm{aug}},x_i^{\mathrm{cf}})\}_{i=1}^{178}.
      \label{eq:srcd-pairs}
\end{align}
The registered schedule shuffles each bank and optimizes on the same
100-element active erase-index set \(\mathcal{I}_{\mathrm{er}}\); the remaining
78 rows stay cached and mapped but do not enter the 200-update optimization.
The two arms share target media, target latents, target-prompt embeddings,
preservation data, initialization, sample order, noise and sigma draws
(hence timesteps), optimizer, and inference settings. Their student embeddings
differ only through factual source identity; the online teacher uses the shared
source-free condition.

\paragraph{Training pressure and falsification.}
Across bindings, source identity varies while the desired source-free outcome
remains fixed. This construction makes identity a less reliable shortcut and
places pressure on the adapter to associate suppression with the registered
source slot and its dependent footprint. It does not guarantee that the model
acquires a causal representation. We test the narrower behavioral
interpretation against the same-backbone Matched Control on all seven
mechanisms, against Generic Paraphrase and frequency-matched Bystander Token
controls on Water and Fracture, and on held-out-source, implicit-footprint,
and specificity cases (\S\ref{subsec:role-evidence}).

\subsection{Counterfactual Distillation and Preservation}
\label{subsec:training-objectives}

\paragraph{Counterfactual erasure update.}
For erasure binding \(i\), let \(z_i^{\mathrm{cf}}\) be the cached latent of
the independently generated source-free target. We sample
\(\epsilon_i\sim\mathcal{N}(0,I)\) and
\(\sigma_i\sim\mathcal{U}(0,1)\), and form
\begin{equation}
  z_{i,\sigma}
  =
  (1-\sigma_i)z_i^{\mathrm{cf}}+\sigma_i\epsilon_i,
  \qquad
  v_i^\star=\epsilon_i-z_i^{\mathrm{cf}}.
  \label{eq:counterfactual-flow-path}
\end{equation}
Here \(z_{i,\sigma}\) is the noised target latent and \(v_i^\star\) is its
flow-matching target. Let
\(\ell\in\{\mathrm{Matched},\mathrm{SRCD}\}\) index the training arm, and let
\(\widehat v_{\theta,\phi_{m,\ell}}(z,t;p)\) denote the prediction of the
frozen base parameters \(\theta\) with the LoRA parameters for mechanism \(m\)
and arm \(\ell\), at model timestep \(t\) under prompt \(p\). The erasure
update is
\begin{align}
  \mathcal{L}_{\mathrm{erase}}^{(i,\ell)}
  ={}&
  \operatorname{MSE}\!\left(
    \widehat v_{\theta,\phi_{m,\ell}}(z_{i,\sigma},t_i;\rho_{i,\ell}),
    v_i^\star
  \right)
  \nonumber\\
  &+
  \lambda_{\mathrm T}
  \operatorname{MSE}\!\left(
    \widehat v_{\theta,\phi_{m,\ell}}(z_{i,\sigma},t_i;\rho_{i,\ell}),
    \operatorname{sg}\!\left[
      \widehat v_{\theta}(z_{i,\sigma},t_i;p_i^{\mathrm{cf}})
    \right]
  \right),
  \qquad \lambda_{\mathrm T}=4,
  \label{eq:counterfactual-erasure-loss}
\end{align}
where \(t_i=1000\sigma_i\), \(\operatorname{sg}\) stops gradients through
the adapter-disabled online teacher, and
\(\rho_{i,\ell}=p_i^{\mathrm{aug}}\) for \methodshort{} but
\(\rho_{i,\ell}=p_i^{\mathrm f}\) for the Matched Control. The
first term fits the flow of the source-free target latent. The second aligns
the factual-conditioned student with the frozen-model prediction under the
source-free target prompt; both predictions receive the same noised latent
and timestep. Only the paired \methodshort{}--Matched
contrast tests an improvement attributable to the source-slot intervention.

\paragraph{Preservation update.}
For generic preservation binding \(j\), let \(z_j^{\mathrm{pr}}\) and
\(p_j^{\mathrm{pr}}\) denote its cached video latent and prompt, and obtain
\(z_{j,\sigma}^{\mathrm{pr}}\) from \(z_j^{\mathrm{pr}}\) by the same
interpolation as in Eq.~\eqref{eq:counterfactual-flow-path}. We match adapted
and frozen predictions under identical inputs:
\begin{equation}
  \mathcal{L}_{\mathrm{preserve}}^{(j,\ell)}
  =
  \lambda_{\mathrm P}
  \operatorname{MSE}\!\left(
    \widehat v_{\theta,\phi_{m,\ell}}
      (z_{j,\sigma}^{\mathrm{pr}},t_j;p_j^{\mathrm{pr}}),
    \operatorname{sg}\!\left[
      \widehat v_{\theta}
        (z_{j,\sigma}^{\mathrm{pr}},t_j;p_j^{\mathrm{pr}})
    \right]
  \right),
  \qquad \lambda_{\mathrm P}=4.
  \label{eq:preservation-loss}
\end{equation}
The student and frozen reference share the latent, sampled noise, timestep,
and prompt. The 36 preservation bindings contain generic non-target scenes,
not specificity cases or hard-negative training examples. This update
constrains drift; the registered specificity, receiver-preservation, quality,
and usability measurements test whether that constraint is respected.

\paragraph{Schedule and scope.}
Training applies Eqs.~\eqref{eq:counterfactual-erasure-loss}
and~\eqref{eq:preservation-loss} on separate, strictly alternating updates;
it does not optimize their sum on a common minibatch. The frozen 200-update
schedule contains 100 erasure and 100 preservation updates, with 100 distinct
erasure bindings selected from the 178-row bank. Only \(\phi_{m,\ell}\) is
optimized, and checkpoint 200 is the sole checkpoint eligible for evaluation.
Matched and \methodshort{} use separate adapters for each mechanism. They share both
objectives, their weights, and all remaining training settings, so the contrast
isolates source-slot replacement rather than the necessity of flow matching,
teacher distillation, preservation matching, or LoRA. Full optimization and
inference settings appear in \appref{app:optimization-details}.
